%% appB --- Pulapki dla ludzi przychodzacych z C#
%%
%% Napisany we wrzesniu 2026, z notes/02-traps.md, ktory jest autorytetem.
%% Kazdy wpis niesie numer z tamtego pliku -- numery nigdy nie sa ponownie
%% uzywane, wiec rozdzial moze sie na wpis powolac -- i nazywa rozdzial,
%% ktory go WYWOLUJE. Wpis, ktorego nie wywoluje zaden rozdzial, jest wada
%% katalogu, a nie faktem o ksiazce; w chwili pisania nie ma takich.

\chapter{Pułapki dla ludzi przychodzących z C\#}
\label{app:traps}

\noindent
Każdy nawyk poniżej kompiluje się w twojej głowie. To właśnie czyni z nich
pułapki, a nie pomyłki: każdy z nich jest poprawnym odruchem co do C\#,
zastosowanym do języka, w którym jest błędny, i większość z nich daje
Pythona, który działa. Książka stawia ci każdy przed nazwaniem go, w
rozdziale, który jest jego właścicielem; ten dodatek jest listą, do której
wracasz.

Dwie rzeczy o numeracji, bo jest nośna. \textbf{Numer nigdy nie jest
używany ponownie}, więc rozdział, który się na niego powołuje, powołuje się
na ten sam wpis na zawsze \dash{} dlatego lista nie jest w kolejności
numerycznej wewnątrz części i dlatego część numerów jest w setkach. Oraz
\textbf{każdy wpis nazywa rozdział, który go wywołuje}, więc każda linia
tutaj jest twierdzeniem o książce, które można sprawdzić, otwierając ten
rozdział.

\section{Część I \dash{} Środowisko uruchomieniowe i narzędzia}
\label{sec:appB-part1}

\trapentry{1}{Wątki są w Pythonie bezużyteczne, więc nie będę sobie nimi
zawracał głowy.}{Blokada serializuje \emph{kod bajtowy}; wątek
zablokowany na I/O nie trzyma blokady. Wątki są właściwym narzędziem do
pracy ograniczonej przez I/O i niewłaściwym do pracy ograniczonej przez
procesor, a eksperyment E1 mierzy oba.
Rozdział~\ref{ch:runtime}.}

\trapentry{2}{Interpreter jest teraz wolnowątkowy, więc blokady już
nie ma.}{Budowa wolnowątkowa to osobny plik binarny, \code{python3.14t},
i wybierasz ją świadomie. Domyślna budowa nadal ma blokadę.
Rozdział~\ref{ch:runtime}.}

\trapentry{3}{Python nie ma kroku kompilacji, więc nie ma tam nic takiego
jak IL.}{Źródło jest kompilowane do kodu bajtowego i cache'owane w
\code{\_\_pycache\_\_}; \api{dis.dis} na funkcji to pokazuje. Nie ma JIT-a,
na którego mógłbyś domyślnie liczyć. Rozdział~\ref{ch:runtime}.}

\trapentry{4}{\code{if \_\_name\_\_ == "\_\_main\_\_"} to
szablonowy kod.}{To jest różnica między modułem, który uruchamia się przy
imporcie, a takim, który uruchamia się przy wykonaniu \dash{} powód, dla
którego listing w tej książce może być i importowalny, i uruchamialny.
Rozdział~\ref{ch:runtime}.}

\trapentry{5}{\code{pip install} kładzie pakiet na maszynie, jak
narzędzie globalne.}{Środowisko jest katalogiem. Instalacja globalna to ta
jedna rzecz, której listingi żadnego z późniejszych rozdziałów nie
przeżyją. Rozdział~\ref{ch:packaging}.}

\trapentry{6}{Aktywuję środowisko, a potem uruchamiam
rzeczy.}{\code{uv run} rozwiązuje środowisko na każde wywołanie. Aktywacja
jest tym nawykiem, który wysyła do CI niewłaściwy interpreter.
Rozdział~\ref{ch:packaging}.}

\trapentry{7}{\code{requirements.txt} jest plikiem blokady.}{To lista
życzeń z zakresami. \code{uv.lock} jest plikiem blokady, jest commitowany,
a \code{uv sync --locked} odmawia od niego odejść.
Rozdział~\ref{ch:packaging}.}

\trapentry{8}{Adnotacja typu jest typem.}{To twierdzenie, którego
środowisko uruchomieniowe nigdy nie sprawdza: parametr z adnotacją
\code{str} może być w czasie działania \code{None} i nic tego nie powie.
Checker jest drugim kompilatorem, którego musisz zaprosić.
Rozdział~\ref{ch:typing}.}

\trapentry{9}{\code{List[int]}, \code{Optional[str]},
\code{Dict[str, Any]}.}{To zapisy z 2019 roku. Na przypiętym
interpreterze jest \code{list[int]}, \code{str | None} i
\code{dict[str, Any]}, a do typów generycznych składnia PEP~695.
Rozdział~\ref{ch:typing}.}

\trapentry{10}{\code{@dataclass} waliduje swoje pola, jak rekord z
konstruktorem.}{Generuje \code{\_\_init\_\_}, \code{\_\_eq\_\_} i
\code{\_\_repr\_\_} i nie sprawdza niczego. Walidacja na granicy jest
zadaniem pydantica. Rozdział~\ref{ch:typing}.}

\trapentry{11}{\code{def f(items=[])} \dash{} opcjonalny parametr
listowy.}{Wartość domyślna jest obliczana raz, przy definicji, i
współdzielona przez każde wywołanie, które jej nie poda. Domyślnie
\code{None} i konstrukcja w środku. Rozdział~\ref{ch:typing}.}

\trapentry{12}{\code{Any} jest jak \code{dynamic}: rzecz, którą mogę
ogrodzić.}{\code{Any} się propaguje. Jedno \code{Any} w łańcuchu zamienia
wszystko dalej w \code{Any}, a checker nie zgłasza zupełnie nic.
Rozdział~\ref{ch:typing}.}

\section{Część II \dash{} Język, odwzorowany}
\label{sec:appB-part2}

\trapentry{13}{Nadpisuję \code{\_\_eq\_\_} i mam równość wartościową, jak
przy nadpisaniu \code{Equals}.}{Zdefiniowanie \code{\_\_eq\_\_} ustawia
\code{\_\_hash\_\_} na \code{None}: obiekt staje się niehaszowalny i nie
może być kluczem \code{dict} ani elementem \code{set}. C\# ostrzega
(CS0659) i jedzie dalej; Python działa. Rozdział~\ref{ch:objects}.}

\trapentry{14}{\code{is} to równość referencji, więc \code{x is 5} jest w
porządku dla małych liczb całkowitych.}{Dwa mechanizmy, a folklor nazywa
tylko jeden. CPython cache'uje liczby całkowite od \code{-5} do
\code{256}; niezależnie od tego kompilator zapisuje równe stałe raz na
obiekt kodu, więc dwa literały \code{257} w jednej funkcji też są jednym
obiektem. Klasyczna demonstracja \emph{nie} odtwarza się więc w skrypcie, a
błąd pojawia się dopiero, gdy wartość przekroczy granicę czasu działania.
Zostaw \code{is} dla \code{None} i wartowników.
Rozdział~\ref{ch:objects}.}

\trapentry{15}{\code{[[]] * n} daje mi \code{n} list.}{Daje jedną listę
\code{n} razy, bo \code{*} kopiuje referencję.
\code{[[] for \_ in range(n)]} daje \code{n}. Rozdział~\ref{ch:objects}.}

\trapentry{16}{Atrybut klasy to pole statyczne.}{Jest współdzielony przez
każdą instancję \emph{i} czytelny przez nie, więc mutowalny atrybut klasy
zmieniony przez \code{self} jest zmieniony dla wszystkich.
\code{@dataclass} odmawia zapisu, który widzi; zapis na zwykłej klasie jest
tym, którego nie widzi. Rozdział~\ref{ch:objects}.}

\trapentry{17}{\code{\_\_str\_\_} to \code{ToString()}.}{\code{\_\_repr\_\_}
jest tym, co pokazuje debugger, REPL i \emph{lista} tych obiektów.
\code{\_\_str\_\_} jest tym, co pokazuje \code{print}, i cofa się do
\code{\_\_repr\_\_}; odwrotnie to nie działa.
Rozdział~\ref{ch:objects}.}

\trapentry{18}{\code{0.1 + 0.2 == 0.3} to błąd Pythona.}{To IEEE~754 w
obu językach, a od .NET Core~3.0 oba nawet wypisują to tak samo, jako
\code{0.30000000000000004}. Używaj \api{math.isclose} do wyliczanych liczb
zmiennoprzecinkowych, a \api{decimal.Decimal} do pieniędzy.
Rozdział~\ref{ch:objects}.}

\trapentry{19}{\code{/} na dwóch liczbach całkowitych to dzielenie
całkowite, jak w C\#.}{\code{/} to dzielenie prawdziwe. \code{//}
zaokrągla w dół, i to ku minus nieskończoności, co pokazuje
\code{-7 // 2 == -4}; odpowiedzią C\# jest w Pythonie \code{int(a / b)}.
Rozdział~\ref{ch:objects}.}

\trapentry{20}{Dekorator to atrybut \dash{} metadane, które czyta
framework.}{Dekorator to funkcja, która uruchamia się w czasie definicji i
może zastąpić to, co dekoruje. To middleware, a nie metadane.
Rozdział~\ref{ch:functions}.}

\trapentry{21}{Wyrażenie listowe to LINQ.}{Wyrażenie listowe jest
zachłanne tam, gdzie \code{IEnumerable} jest odroczony, a wyrażenie
generatorowe jest odroczone \emph{i jednoprzebiegowe}, podczas gdy zapytanie
LINQ uruchamia się ponownie przy każdym przejściu. Materializuj raz, celowo.
Rozdział~\ref{ch:functions}.}

\trapentry{22}{Domknięcie w pętli przechwytuje wartość zmiennej
pętli.}{Przechwytuje zmienną, wiązaną późno, więc każde domknięcie widzi
ostatnią wartość. Zwiąż argumentem domyślnym albo
\api{functools.partial}. Ruff łapie to jako B023, co czyni z tego jedyną
pułapkę w tym rozdziale, którą narzędziówka znajduje za ciebie.
Rozdział~\ref{ch:functions}.}

\trapentry{23}{\code{s += piece} w pętli jest w porządku; napisy to
napisy.}{Na CPythonie zwykła pętla jest \emph{liniowa}: interpreter ma
specjalizację, która zmienia rozmiar w miejscu, kiedy lewy argument nie ma
innej referencji. Jest znów kwadratowa w chwili, gdy cokolwiek innego trzyma
referencję \dash{} lista, domknięcie, atrybut \dash{} bez żadnej zmiany w
miejscu wywołania. \code{"".join(pieces)} nie ma złego przypadku.
Rozdział~\ref{ch:functions}.}

\trapentry{24}{Generator to \code{yield return}.}{Jest nim, łącznie z
\code{finally}, które wykonuje się przy wcześniejszym wyjściu, czyli z
\code{IEnumerator.Dispose}. Ma też \code{send()}, \code{close()} i
wynoszoną wartość zwracaną \dash{} korutynę, zanim istniało \code{async},
dlatego \code{Generator} ma trzy parametry typu tam, gdzie
\code{IEnumerable} ma jeden. Rozdział~\ref{ch:functions}.}

\trapentry{140}{\code{case ACTIVE:} porównuje ze stałą, jak robi to ramię
\code{switch}.}{Goła nazwa we wzorcu \textbf{przechwytuje}: wiąże każdy
podmiot z \code{ACTIVE} i ramię zawsze pasuje. Kompilator odmawia, gdy
następuje po nim kolejny przypadek, więc kompiluje się w ciszy tylko jako
ostatni \dash{} czyli dokładnie tam, gdzie przeżywa przegląd. Użyj nazwy z
kropką albo strażnika. Rozdział~\ref{ch:functions}.}

\trapentry{25}{\code{except:} łapie wszystko, jak \code{catch
\{\}}.}{Łapie też \code{KeyboardInterrupt} i \code{SystemExit}, dlatego
procesu nie da się zatrzymać. Najwyżej \code{except Exception:}, i loguj
ślad stosu. Ruff zgłasza to jako E722.
Rozdział~\ref{ch:errors}.}

\trapentry{26}{Łapię \code{Exception}, loguję i jadę dalej \dash{} to
zapobiegawcze.}{Zamienia awarię w cichą błędną odpowiedź. EAFP oznacza
łapanie wyjątku, którego się spodziewasz, a nie każdego. To jedyna pułapka
w tym rozdziale, której nic w narzędziówce książki nie zgłasza, i dlatego
właśnie ta dociera na produkcję. Rozdział~\ref{ch:errors}.}

\trapentry{27}{\code{raise e} wewnątrz \code{except} rzuca ponownie, jak
\code{throw ex;}.}{Python trzyma ślad stosu na \emph{obiekcie} wyjątku,
więc \code{raise e} niczego nie ucina: każda ramka pod nim przeżywa, a
linia ponownego rzucenia jest \emph{dodawana}. Gołe \code{raise} jest nadal
lepszym nawykiem, bo zdublowana ramka to szum, ale reguła z C\# się nie
przenosi. Co ponowne rzucenie może stracić, to powiązanie: bez
\code{from} \code{\_\_cause\_\_} zostaje \code{None}.
Rozdział~\ref{ch:errors}.}

\trapentry{28}{\code{return} w \code{finally} jest nieszkodliwy.}{Połyka
każdy wyjątek w locie i wraca, jakby nic się nie stało. Nie po cichu na
przypiętym interpreterze \dash{} PEP~765 każe kompilatorowi wyemitować
\code{SyntaxWarning}, a ruff zgłasza B012 i SIM107 \dash{} ale nadal połyka,
bo ostrzeżenie nie jest błędem. Rozdział~\ref{ch:errors}.}

\trapentry{29}{Adnotacja metody mówi mi, co ona rzuca.}{Nic w systemie
typów nie niesie wyjątków. Docstring niesie, i nic tego nie weryfikuje.
Python nie ma wyjątków kontrolowanych \dash{} ani C\# ich nie ma, więc źle
przenosi się narzędziówka wokół nich, a nie język.
Rozdział~\ref{ch:errors}.}

\trapentry{30}{\code{KeyError} znaczy, że coś poszło
źle.}{\code{KeyError}, \code{StopIteration} i \code{AttributeError} są
protokołem: odczyt ze \code{dict}, koniec iteratora i \api{getattr} mówią
przez nie. Rozdział~\ref{ch:errors}.}

\trapentry{31}{Prawdziwościowość to \code{bool}, jak w C\#.}{Puste
kontenery, zero, \code{None} i puste napisy są fałszem. \code{if items:}
jest idiomatyczne, a \code{if items is not None:} zadaje inne pytanie.
Rozdział~\ref{ch:errors}.}

\trapentry{150}{Klasa wyjątku to \code{FooException}.}{Przyrostkiem
Pythona jest \code{Error}, a N818 w ruffie zgłasza klasę bez niego:
\code{JobUnavailable} oblewa lint, dopóki nie jest
\code{JobUnavailableError}. Rozdział~\ref{ch:errors}.}

\trapentry{32}{Moduł to przestrzeń nazw, więc import jest darmowy i
czysty.}{Moduł to obiekt, który wykonuje się raz, od góry do dołu. Efekt
uboczny przy imporcie wykonuje się dla każdego importującego, a import
cykliczny to dwa moduły w połowie wykonane.
Rozdział~\ref{ch:imports}.}

\trapentry{33}{\code{from x import *} to \code{using x;}.}{Kopiuje każdą
publiczną nazwę do modułu importującego i ukrywa, skąd cokolwiek przyszło.
Używaj \code{import x} albo \code{from x import name}.
Rozdział~\ref{ch:imports}.}

\trapentry{34}{Mogę nazwać zmienną \code{list}, \code{id} albo
\code{type}.}{Przesłania to wbudowaną nazwę do końca zasięgu, a awaria
przychodzi trzy funkcje później. Rozdział~\ref{ch:imports}.}

\trapentry{35}{Potrzebuję kontenera DI.}{Korzeń kompozycji jest funkcją.
\api{functools.partial} i \code{Protocol} robią to, co robił kontener, a
\code{Depends} z FastAPI jest tym jednym kontenerem, który większość
czytelników spotka. Rozdział~\ref{ch:imports}.}

\section{Część III \dash{} Współbieżność}
\label{sec:appB-part3}

\trapentry{36}{Korutyna to \code{Task}: wywołanie jej ją
uruchamia.}{Korutyna jest zimna. Wywołanie jej buduje obiekt, który nic nie
robi, dopóki nie jest oczekiwany albo zaplanowany, a zapomniany
\code{await} jest ostrzeżeniem, a nie błędem.
Rozdział~\ref{ch:asyncio}.}

\trapentry{37}{\code{create\_task} uruchamia zadanie
natychmiast.}{Planuje je; ciało nie działa, dopóki wołający nie ustąpi.
Rozdział~\ref{ch:asyncio}.}

\trapentry{38}{\code{.Result} na zadaniu powoduje zakleszczenie, więc
używam \code{ConfigureAwait(false)}.}{Nie ma kontekstu synchronizacji ani
\code{ConfigureAwait}. Jest jedna pętla, a jedno blokujące wywołanie w
niej zamraża każdą inną korutynę, bez błędu i bez linii w logu.
Rozdział~\ref{ch:asyncio}.}

\trapentry{39}{\code{gather} to \code{WhenAll}.}{\code{gather} osierocą
swoje rodzeństwo, gdy jedno padnie; \code{TaskGroup} je anuluje. Oba są
identyczne co do szybkości, więc wybór dotyczy wyłącznie semantyki awarii.
Rozdział~\ref{ch:asyncio}.}

\trapentry{40}{Anulowanie to token, który odpytuję, więc łapanie
\code{Exception} jest tym, jak je tracę.}{Na opak w obu połowach. To
\code{CancelledError}, dostarczany przy \code{await}, i dziedziczy po
\code{BaseException} \dash{} więc \code{except Exception} jest klauzulą,
która je \emph{przepuszcza}. Połyka je gołe \code{except:},
\code{except BaseException:} albo blok \code{except asyncio.CancelledError}
bez \code{raise} pod nim. Rozdział~\ref{ch:asyncio}.}

\trapentry{41}{\code{HttpClient} jest singletonem, więc
\code{httpx.AsyncClient} też.}{Rada co do czasu życia się przenosi
\dash{} skonstruuj raz, współdziel, zamknij \dash{} a wartości domyślne
nie: domyślny klient ma już pięciosekundowy termin, nakładany osobno na
połączenie, odczyt, zapis i pulę, a nie na żądanie jako całość.
Rozdział~\ref{ch:asyncio}.}

\trapentry{170}{Na \code{Task} można czekać dwa razy, więc na korutynę
też.}{\code{Task} to uchwyt do pracy już trwającej i rozdaje swój wynik
tyle razy, ile poprosisz. Korutyna \emph{jest} tą pracą, a oczekiwanie na
nią po raz drugi podnosi \code{RuntimeError}.
Rozdział~\ref{ch:asyncio}.}

\section{Część IV \dash{} Wdrażanie}
\label{sec:appB-part4}

\trapentry{42}{\code{Depends} to kontener DI, więc zależność jest
singletonem.}{Domyślnie jest na żądanie. Zależność z \code{yield} to
zasięg z zamknięciem, a stan o czasie życia aplikacji idzie do lifespanu.
Rozdział~\ref{ch:services}.}

\trapentry{43}{Niepowodzenie walidacji to 400 z
\code{ProblemDetails}.}{To 422 z listą błędów pydantica, a kształt należy
do pydantica, a nie do frameworka. Rozdział~\ref{ch:services}.}

\trapentry{44}{Ustawienia biorą się z \code{IOptions<T>} i
\code{appsettings.json}.}{\pkg{pydantic-settings} czyta środowisko, plik
\code{.env} i katalog sekretów do zwalidowanego modelu. Nie ma łańcucha
dostawców JSON. Rozdział~\ref{ch:services}.}

\trapentry{45}{Więcej workerów uvicorna to większa
przepustowość.}{Workery to procesy: każdy ma własną pętlę, własną pulę i
własną pamięć. E5 to zmierzył i dla oczekiwanego upstreamu cztery workery
nie ruszyły przepustowości, mnożąc przy tym pamięć.
Rozdział~\ref{ch:services}.}

\trapentry{180}{Middleware działa w kolejności, w jakiej je
zarejestrowałem, więc pierwszy jest tym zewnętrznym.}{Ostatni
zarejestrowany jest najbardziej zewnętrzny: \code{add\_middleware} ze
starlette wstawia na początek stosu. Nic nie ostrzega, a obie kolejności
różnią się tylko wtedy, gdy coś zależy od sekwencji.
Rozdział~\ref{ch:services}.}

\trapentry{181}{\code{Field(alias=...)} to \code{JsonPropertyName}:
zmienia nazwę pola na drucie.}{Zmienia ją także na wejściu, więc modelu
nie da się już skonstruować po nazwie pola \dash{} a checker typów mówi to
wcześniej niż test. \code{serialization\_alias} jest tym tylko
wyjściowym. Rozdział~\ref{ch:services}.}

\trapentry{182}{\code{--workers 4} rozkłada mój ruch na cztery
procesy.}{Workery współdzielą jedno gniazdo nasłuchujące, a ten worker,
który pierwszy jest w \code{accept()}, trzyma połączenie przez całe jego
życie, więc równoważenie dzieje się na \emph{połączenie}, a nie na
żądanie. Rozdział~\ref{ch:services}.}

\trapentry{46}{\code{Session} to \code{DbContext}, więc robię
\code{commit()} i obiekty zostają używalne.}{\code{expire\_on\_commit} jest
domyślnie prawdą, więc każdy atrybut jest przeładowywany przy następnym
dostępie \dash{} czyli zapytanie w pętli, której nie napisałeś.
Rozdział~\ref{ch:data}.}

\trapentry{47}{Relacje ładują się, kiedy je czytam, jak w EF.}{Leniwie
domyślnie w obu, pod sesją \emph{synchroniczną}, i N+1 jest identyczne;
\code{selectinload} to \code{Include}. Pod sesją asynchroniczną ten sam
odczyt podnosi \code{MissingGreenlet} zamiast ładować, więc błąd jest tam
głośny, a tu cichy. Rozdział~\ref{ch:data}.}

\trapentry{48}{Autogenerate Alembica jest migracją.}{To diff modeli
względem bazy i gubi zmiany nazw, zmiany typów i więzy. Zmiana nazwy
wychodzi jako dodanie plus usunięcie, co gubi dane. Czytaj każdą
wygenerowaną migrację. Rozdział~\ref{ch:data}.}

\trapentry{49}{Potrzebuję repozytorium nad ORM-em.}{\code{Session} jest
jednostką pracy, a wzorzec repozytorium ją dubluje.
Rozdział~\ref{ch:data} mówi, dlaczego to zwykle błąd w tym miejscu.}

\trapentry{50}{Klasa testowa na fixture, jak
\code{IClassFixture}.}{Fixture'y pytesta to funkcje z zasięgiem. Klasa jest
opcjonalna i zwykle jej nie ma. Rozdział~\ref{ch:testing}.}

\trapentry{51}{Mockuję tam, gdzie rzecz jest zdefiniowana.}{Podmieniaj
tam, gdzie nazwa jest \emph{wyszukiwana} \dash{} w module importującym
\dash{} albo podmiana nic nie robi, a test przechodzi z niewłaściwego
powodu. To, którą nazwę rozwiązuje wywołanie, zdecydował import, więc
\code{import x} plus \code{x.y()} naprawdę podmienia się w \code{x.y}.
Rozdział~\ref{ch:testing}.}

\trapentry{52}{\code{assert} jest do buildów debugowych.}{pytest
przepisuje \code{assert}, żeby sam się tłumaczył; to jest biblioteka
asercji. Połowa nawyku jest słuszna: pod \code{-O} \code{assert} w
\emph{aplikacji} naprawdę znika, więc jest do testów i niezmienników, a
nigdy do walidacji wejścia. Rozdział~\ref{ch:testing}.}

\trapentry{200}{Raport z testów to raport z testów; ta sama awaria
wypisuje wszędzie to samo.}{pytest raportuje na serwerze build\-owym
inaczej, i to celowo: kilka miejsc pyta, czy ustawione jest \code{CI} albo
\code{BUILD\_NUMBER}, i wtedy różnica sekwencji jest wypisywana w całości, a
długie wyjście nie jest ucinane. Transkrypt, plik wzorcowy albo zrzut
ekranu z awarii jest twierdzeniem o jednej z dwóch maszyn.
Rozdział~\ref{ch:testing}.}

\trapentry{53}{\api{logging.basicConfig} i gotowe.}{Konfiguracja modułu
standardowego jest globalna i wrażliwa na kolejność importów, a drugie
wywołanie jest cichym brakiem działania. To jest powód, dla którego istnieje
structlog. Rozdział~\ref{ch:observability}.}

\trapentry{54}{Identyfikator korelacji w polu statycznym, jak
\code{AsyncLocal}.}{\api{contextvars} to \code{AsyncLocal}. Zmienna na
poziomie modułu jest współdzielona przez każde żądanie na pętli.
Rozdział~\ref{ch:observability}.}

\trapentry{55}{\code{FROM python} i \code{pip install} w
Dockerfile'u.}{Wieloetapowo, z \code{uv sync --frozen --no-dev},
użytkownikiem bez roota i dwiema zmiennymi środowiskowymi.
Rozdział~\ref{ch:observability}.}

\section{Część V \dash{} Python do pracy z AI}
\label{sec:appB-part5}

\trapentry{56}{SDK jest magią.}{Każdy SDK do AI to model pydantica,
klient httpx i iterator strumieniowy. Przeczytaj zainstalowany, przypięty
pakiet. Rozdział~\ref{ch:aitoolkit}.}

\trapentry{57}{\code{model\_validate\_json} to
\code{JsonSerializer.Deserialize}.}{Tryb luźny konwertuje tam, gdzie
\code{System.Text.Json} odmawia; tryb ścisły jest tym, czego oczekuje
inżynier C\#, a E8 wycenia oba \dash{} na nic.
Rozdział~\ref{ch:aitoolkit}.}

\trapentry{58}{Notatnik jest tam, gdzie dzieje się Python.}{Notatnik to
REPL z pamięcią każdej komórki, którą uruchomiłeś w dowolnej kolejności.
Nic w tej książce nim nie jest, a rozdział~\ref{ch:aitoolkit} mówi, kiedy
jest właściwy.}

\trapentry{220}{SDK używa \pkg{httpx}, który przypiąłem.}{Oba SDK zależą
od dystrybucji \pkg{httpx2}, a nie od \pkg{httpx}, więc sprawdzenie
\api{isinstance} względem tej, którą przypiąłeś, daje \code{False} na
kliencie, który jest oczywiście klientem httpx. Z dwóch SDK jeden odmawia
niedopasowanego klienta, a drugi go przyjmuje.
Rozdział~\ref{ch:aitoolkit}.}

\trapentry{59}{Asercja na śladzie potrzebuje modelu, żeby ją
ocenić.}{Pierwsza warstwa instrumentu, który ta książka portuje, jest
deterministyczna: działa bez modelu, i właśnie dlatego może działać w CI.
Rozdział~\ref{ch:traceassert}.}
